% Options for packages loaded elsewhere
\PassOptionsToPackage{unicode}{hyperref}
\PassOptionsToPackage{hyphens}{url}
\PassOptionsToPackage{dvipsnames,svgnames,x11names}{xcolor}
\documentclass[
  9pt,
]{article}
\usepackage{xcolor}
\usepackage[margin=1.5cm]{geometry}
\usepackage{amsmath,amssymb}
\setcounter{secnumdepth}{-\maxdimen} % remove section numbering
\usepackage{iftex}
\ifPDFTeX
  \usepackage[T1]{fontenc}
  \usepackage[utf8]{inputenc}
  \usepackage{textcomp} % provide euro and other symbols
\else % if luatex or xetex
  \usepackage{unicode-math} % this also loads fontspec
  \defaultfontfeatures{Scale=MatchLowercase}
  \defaultfontfeatures[\rmfamily]{Ligatures=TeX,Scale=1}
\fi
\usepackage{lmodern}
\ifPDFTeX\else
  % xetex/luatex font selection
\fi
% Use upquote if available, for straight quotes in verbatim environments
\IfFileExists{upquote.sty}{\usepackage{upquote}}{}
\IfFileExists{microtype.sty}{% use microtype if available
  \usepackage[]{microtype}
  \UseMicrotypeSet[protrusion]{basicmath} % disable protrusion for tt fonts
}{}
\makeatletter
\@ifundefined{KOMAClassName}{% if non-KOMA class
  \IfFileExists{parskip.sty}{%
    \usepackage{parskip}
  }{% else
    \setlength{\parindent}{0pt}
    \setlength{\parskip}{6pt plus 2pt minus 1pt}}
}{% if KOMA class
  \KOMAoptions{parskip=half}}
\makeatother
\usepackage{longtable,booktabs,array}
\usepackage{calc} % for calculating minipage widths
% Correct order of tables after \paragraph or \subparagraph
\usepackage{etoolbox}
\makeatletter
\patchcmd\longtable{\par}{\if@noskipsec\mbox{}\fi\par}{}{}
\makeatother
% Allow footnotes in longtable head/foot
\IfFileExists{footnotehyper.sty}{\usepackage{footnotehyper}}{\usepackage{footnote}}
\makesavenoteenv{longtable}
\setlength{\emergencystretch}{3em} % prevent overfull lines
\providecommand{\tightlist}{%
  \setlength{\itemsep}{0pt}\setlength{\parskip}{0pt}}
\usepackage{amssymb}
\usepackage{bookmark}
\IfFileExists{xurl.sty}{\usepackage{xurl}}{} % add URL line breaks if available
\urlstyle{same}
\hypersetup{
  colorlinks=true,
  linkcolor={blue},
  filecolor={Maroon},
  citecolor={Blue},
  urlcolor={blue},
  pdfcreator={LaTeX via pandoc}}

\author{}
\date{}

\begin{document}

\subsection{1. Dataset Overview}\label{dataset-overview}

\begin{itemize}
\tightlist
\item
  \textbf{Total records:} 1,297,015
\item
  \textbf{Columns:} 35
\item
  \textbf{Sites:} 106
\end{itemize}

\begin{longtable}[]{@{}
  >{\raggedright\arraybackslash}p{(\linewidth - 10\tabcolsep) * \real{0.1188}}
  >{\raggedleft\arraybackslash}p{(\linewidth - 10\tabcolsep) * \real{0.0891}}
  >{\raggedleft\arraybackslash}p{(\linewidth - 10\tabcolsep) * \real{0.0891}}
  >{\raggedright\arraybackslash}p{(\linewidth - 10\tabcolsep) * \real{0.1683}}
  >{\raggedright\arraybackslash}p{(\linewidth - 10\tabcolsep) * \real{0.2673}}
  >{\raggedright\arraybackslash}p{(\linewidth - 10\tabcolsep) * \real{0.2673}}@{}}
\toprule\noalign{}
\begin{minipage}[b]{\linewidth}\raggedright
dataset
\end{minipage} & \begin{minipage}[b]{\linewidth}\raggedleft
rows
\end{minipage} & \begin{minipage}[b]{\linewidth}\raggedleft
sites
\end{minipage} & \begin{minipage}[b]{\linewidth}\raggedright
sampling
\end{minipage} & \begin{minipage}[b]{\linewidth}\raggedright
first\_\allowbreak{}ts
\end{minipage} & \begin{minipage}[b]{\linewidth}\raggedright
last\_\allowbreak{}ts
\end{minipage} \\
\midrule\noalign{}
\endhead
\bottomrule\noalign{}
\endlastfoot
goes\_\allowbreak{}pvdaq & 88806 & 8 & 15 minutes & 2019-01-01 00:00:00+00:00 &
2019-09-30 23:45:00+00:00 \\
uk\_\allowbreak{}pv & 1208209 & 98 & 30 minutes & 2019-01-01 08:00:00+00:00 &
2020-12-31 16:00:00+00:00 \\
\end{longtable}

\subsubsection{1.1 Organization}\label{organization}

Two source datasets (\texttt{uk\_\allowbreak{}pv} satellite HRV crops,
\texttt{goes\_\allowbreak{}pvdaq} GOES-16 crops) merged into a single table
\texttt{dataset\_\allowbreak{}all.parquet} --- one row per
\texttt{(site\_\allowbreak{}id,\ timestamp\_\allowbreak{}utc)}. Each row pairs a PV power reading
and its weather/solar covariates with a co-registered image frame;
pixels live in \texttt{images\_\allowbreak{}all.h5}, referenced by the integer
\texttt{image\_\allowbreak{}*\_\allowbreak{}index} pointers. \texttt{norm\_\allowbreak{}power} (power \(\div\)
installed capacity) is the model target; splits are cross-plant on
\texttt{site\_\allowbreak{}id}.

\subsubsection{1.2 Columns}\label{columns}

\begin{longtable}[]{@{}
  >{\raggedright\arraybackslash}p{(\linewidth - 4\tabcolsep) * \real{0.3333}}
  >{\raggedright\arraybackslash}p{(\linewidth - 4\tabcolsep) * \real{0.3333}}
  >{\raggedright\arraybackslash}p{(\linewidth - 4\tabcolsep) * \real{0.3333}}@{}}
\toprule\noalign{}
\begin{minipage}[b]{\linewidth}\raggedright
Column
\end{minipage} & \begin{minipage}[b]{\linewidth}\raggedright
Meaning
\end{minipage} & \begin{minipage}[b]{\linewidth}\raggedright
Role
\end{minipage} \\
\midrule\noalign{}
\endhead
\bottomrule\noalign{}
\endlastfoot
\texttt{dataset} & Source dataset: \texttt{uk\_\allowbreak{}pv} or
\texttt{goes\_\allowbreak{}pvdaq} & id --- cross-region split key \\
\texttt{site\_\allowbreak{}id} & PV plant identifier & id --- cross-plant split
key \\
\texttt{timestamp\_\allowbreak{}utc} & UTC sample time & index --- windowing \\
\texttt{power\_\allowbreak{}w} & Measured PV power output (W) & raw target; not fed
directly \\
\texttt{installed\_\allowbreak{}power\_\allowbreak{}w} & Nameplate capacity (W) & normalizer \\
\texttt{norm\_\allowbreak{}power} & power\_\allowbreak{}w \(\div\) installed\_\allowbreak{}power\_\allowbreak{}w &
\textbf{target} \\
\texttt{temperature\_\allowbreak{}2m} & 2 m air temperature (\(^\circ\)C) & covariate ---
observed \\
\texttt{shortwave\_\allowbreak{}radiation} & Surface downwelling shortwave (W/m\textsuperscript{2}) &
covariate --- observed \\
\texttt{direct\_\allowbreak{}radiation} & Direct-beam component (W/m\textsuperscript{2}) & covariate
--- observed \\
\texttt{diffuse\_\allowbreak{}radiation} & Diffuse component (W/m\textsuperscript{2}) & covariate ---
observed \\
\texttt{direct\_\allowbreak{}normal\_\allowbreak{}irradiance} & DNI (W/m\textsuperscript{2}) & covariate ---
observed \\
\texttt{cloudcover} & Total cloud cover (\%) & covariate --- observed \\
\texttt{windspeed\_\allowbreak{}10m} & 10 m wind speed (m/s) & covariate ---
observed \\
\texttt{precipitation} & Precipitation (mm) & covariate --- observed \\
\texttt{solar\_\allowbreak{}zenith} & Solar zenith angle (deg) & covariate ---
deterministic \\
\texttt{solar\_\allowbreak{}azimuth} & Solar azimuth angle (deg) & covariate ---
deterministic \\
\texttt{clearsky\_\allowbreak{}ghi} & Clear-sky GHI model (W/m\textsuperscript{2}) & covariate ---
deterministic \\
\texttt{doy\_\allowbreak{}sin}, \texttt{doy\_\allowbreak{}cos} & Cyclic day-of-year encoding &
covariate --- deterministic \\
\texttt{solar\_\allowbreak{}time} & Local apparent solar time (hours) & covariate ---
deterministic \\
\texttt{image\_\allowbreak{}h5\_\allowbreak{}index} & Row index into \texttt{images\_\allowbreak{}all.h5} &
vision input pointer \\
\texttt{latitude}, \texttt{longitude} & Site coordinates & metadata ---
unused (can be used for graph building) \\
\texttt{station\_\allowbreak{}id} & Weather station tag & metadata --- unused \\
\texttt{camera\_\allowbreak{}id} & Imagery channel (HRV/GOES) & metadata ---
unused \\
\texttt{outage\_\allowbreak{}flag} & Flag: outage / zero-output anomaly & curation
only, unused \\
\texttt{night\_\allowbreak{}clamped} & Flag: nighttime power clamped to 0 & curation
only, unused \\
\end{longtable}

\newpage

\subsubsection{1.3 Construction \& sources}\label{construction-sources}

The dataset is assembled fusing four public sources into one
standardized \texttt{(site\_\allowbreak{}id,\ timestamp\_\allowbreak{}utc)} table plus
co-registered image frames:

\begin{itemize}
\tightlist
\item
  \textbf{PV power --- UK}
  (\href{https://huggingface.co/datasets/openclimatefix/uk_pv}{\texttt{openclimatefix/uk\_\allowbreak{}pv}},
  HF): 30-minutely generation (2019--2020) and per-site metadata
  (rounded lat/lon, \texttt{kWp} capacity). \texttt{generation\_\allowbreak{}Wh} over
  30 min is converted to watts (\(\times\)2); short gaps (\(\leq\) 3 steps) are linearly
  interpolated, unfixable ones dropped.
\item
  \textbf{UK satellite imagery} --- EUMETSAT \textbf{SEVIRI RSS HRV}
  from the Open Climate Fix public GCS bucket
  (\texttt{gs://public-datasets-eumetsat-solar-forecasting/satellite/EUMETSAT/SEVIRI\_\allowbreak{}RSS/v4/\textless{}year\textgreater{}\_\allowbreak{}hrv.zarr}),
  \textasciitilde1 km/px. Frames are reprojected (pyproj) to each site
  and cropped to \textbf{128\(\times\)128 px (\textasciitilde128 km)} for
  cloud-advection context.
\item
  \textbf{PV power + imagery --- US} (\texttt{goes\_\allowbreak{}pvdaq}): NREL
  \textbf{PVDAQ} systems with site metadata (lat/lon, nameplate power)
  from the \href{https://oedi-data-lake.s3.amazonaws.com/pvdaq/}{OEDI
  data lake}, paired with \textbf{GOES-16} satellite crops (10 sites).
\item
  \textbf{Weather covariates} ---
  \href{https://open-meteo.com/en/docs/historical-weather-api}{Open-Meteo
  Archive API}: 8 hourly variables (2 m temperature,
  shortwave/direct/diffuse radiation, DNI, cloud cover, 10 m wind speed,
  precipitation), fetched per site and joined to each row by
  nearest-timestamp (\texttt{merge\_\allowbreak{}asof}).
\end{itemize}

\textbf{Operations.} Per frame: NaN \(\rightarrow\) 0, clip to {[}0,1{]}, keep only
non-empty crops; images are written as per-site PNGs / a single
\texttt{uint8} HDF5 (millions of small files are impractical). Numeric
rows from both sources are concatenated, enriched with
\texttt{installed\_\allowbreak{}power\_\allowbreak{}w} + coordinates, then joined with the weather
track and saved.

\subsubsection{1.4 Other public dataset --- ClimateHackAI
2023}\label{other-public-dataset-climatehackai-2023}

The \textbf{ClimateHackAI 2023} competition dataset
(\href{https://doxaai.com/competition/climatehackai-2023/overview}{DOXA}
\(\cdot\)
\href{https://huggingface.co/datasets/climatehackai/climatehackai-2023}{HF},
\textasciitilde600 GB) is a multimodal PV-nowcasting corpus over
\textbf{Great Britain}. Per site it provides \textbf{PV generation} as a
fraction of installed capacity at \textbf{5-min} resolution, with
metadata (latitude, longitude, orientation, tilt, installed capacity);
two \textbf{EUMETSAT} satellite streams at 5-min cadence ---
\textbf{HRV} \texttt{{[}12,128,128{]}} and \textbf{non-HRV} 11-channel
\texttt{{[}12,128,128,11{]}} crops centred on each site; \textbf{DWD
ICON-EU} numerical weather forecasts (38 variables, T\(-\)1h\ldots T+4h,
128\(\times\)128 grids); and \textbf{ECMWF CAMS} air-quality forecasts (13
variables \(\times\) 8 altitudes). The competition task: from \textbf{1 h of
history} (PV + satellite) plus the hourly weather/air-quality forecast
steps, predict the \textbf{next 4 h} of PV output --- 48 values at 5-min
intervals --- scored by mean absolute error (target \textless{} 0.15).

It is structured differently from the dataset of record: built for
\textbf{nowcasting}, it is organised as short fixed windows (1 h in \(\rightarrow\) 4
h out) at 5-min resolution with per-window satellite/NWP/air-quality
tensors, rather than a single long tall table of
\texttt{(site,\ timestamp)} rows spanning history and horizon. That
short-window, high-cadence, forecast-tensor layout targets
minutes-to-hours nowcasting and would need re-windowing and re-sampling
before it could serve the longer-history, cross-plant protocol used
here.

Note: this is a \textbf{competition dataset with no accompanying
published paper}; it is documented only through the competition overview
and Hugging Face repository, so there is no peer-reviewed reference to
cite for its construction or provenance.

\subsubsection{1.5 Other public dataset --- MMSP
(FusionSF)}\label{other-public-dataset-mmsp-fusionsf}

The \textbf{MMSP} (Multi-Modal Solar Power) dataset is released with
\textbf{FusionSF}
(\href{https://arxiv.org/abs/2402.05823}{arXiv:2402.05823} \(\cdot\)
\href{https://github.com/MAZiqing/FusionSF}{code} \(\cdot\)
\href{https://drive.google.com/drive/folders/1qGVOw-hAVQlO3n-1d4ZNHvL42L9PkdBK}{data}),
a multimodal PV corpus over \textbf{88 geographically dispersed solar
plants across a Chinese province (\textasciitilde157,100 km\textsuperscript{2})}, spanning
\textbf{Jan 2021 -- Jun 2022} at \textbf{hourly} resolution (downsampled
from 10-min). Each record fuses three modalities: \textbf{historical PV
power}; \textbf{Himawari-8/9 satellite imagery} cropped to \textbf{64\(\times\)64
px} (from 640\(\times\)640) --- 1 channel in the small variant \textbf{MMSP(S)}
(10 plants), 4 visible/near-infrared channels in the full
\textbf{MMSP(L)} (88 plants); and \textbf{ECMWF NWP} (17 weather
features: radiation, temperature, pressure, cloud cover, wind). The task
is \textbf{day-ahead} forecasting: 24 h of history \(\rightarrow\) next 24 h. (The
paper reports the method deployed on 300+ plants / \textgreater15 GW in
production; MMSP is the released research subset.)

Like ClimateHackAI, it differs structurally from the dataset of record
--- hourly cadence, fixed day-ahead 24-in \(\rightarrow\) 24-out windows over China
with Himawari imagery and per-window NWP tensors --- so it would need
re-windowing/re-sampling before serving the longer-history, cross-plant
protocol used here.

\newpage

\subsection{2. Data Splits}\label{data-splits}

Cross-plant, seeded per-dataset partition over the 106 sites. Counts
combine both datasets.

\begin{longtable}[]{@{}lcc@{}}
\toprule\noalign{}
& Plants & Rows \\
\midrule\noalign{}
\endhead
\bottomrule\noalign{}
\endlastfoot
\textbf{Train} & 75 & 918 158 \\
\textbf{Validation} & 16 & 193 553 \\
\textbf{Test} & 15 & 185 304 \\
\end{longtable}

Per-dataset plants --- uk\_\allowbreak{}pv: 69/15/14 (train/val/test); goes\_\allowbreak{}pvdaq:
6/1/1.

\subsection{3. Configurations}\label{configurations}

\begin{longtable}[]{@{}
  >{\raggedright\arraybackslash}p{(\linewidth - 2\tabcolsep) * \real{0.4444}}
  >{\centering\arraybackslash}p{(\linewidth - 2\tabcolsep) * \real{0.5556}}@{}}
\toprule\noalign{}
\begin{minipage}[b]{\linewidth}\raggedright
Dimension
\end{minipage} & \begin{minipage}[b]{\linewidth}\centering
Specs
\end{minipage} \\
\midrule\noalign{}
\endhead
\bottomrule\noalign{}
\endlastfoot
\textbf{History} & 14 days (steps: 672 uk and 1344 pvdaq) \\
\textbf{Horizon} & 6 hours (steps: 12 uk and 24 pvdaq) \\
\textbf{Frames} & 8 frames @ 60-min cadence (last 8 h ending at forecast
time) \\
\end{longtable}

Models should be tested also at intra-hour frame cadence. On
\texttt{uk\_\allowbreak{}pv} 30 min and on \texttt{goes\_\allowbreak{}pvdaq} 15 min can be used
(longer training).

\subsection{2. Leaderboard}\label{leaderboard}

\textbf{Rank} is the global position by skill score (SS \(\uparrow\)) across all
models; each subtable is sorted best-first. Metrics: SS \(\uparrow\), R\textsuperscript{2} \(\uparrow\), CRPS \(\downarrow\).
Reference floor is \texttt{smart\_\allowbreak{}persistence} (SS \(\equiv\) 0).

\textbf{Metric intuition.} All point metrics run on the
capacity-normalized target (\texttt{norm\_\allowbreak{}power} \(\in\) {[}0,1{]}) over
daylight, valid steps, so values read as fraction-of-nameplate error and
compare across plants.

\begin{itemize}
\tightlist
\item
  \textbf{R\textsuperscript{2}} \(\uparrow\) --- squared Pearson correlation between predicted and
  true power, r\textsuperscript{2} = corr(true, pred)\textsuperscript{2}, pooled across test plants on
  daylight, valid steps. Measures how much of the \emph{variance} the
  forecast tracks (shape/timing), independent of bias or scale. 1 =
  perfectly correlated, 0 = no linear relation.
\item
  \textbf{SS (skill score)} \(\uparrow\) --- how much a model beats the
  \texttt{smart\_\allowbreak{}persistence} baseline on root-mean-square error: 0 = no
  better than baseline, 1 = perfect, negative = worse.
\item
  \textbf{CRPS} \(\downarrow\) --- scores the \emph{whole predictive distribution}
  (via quantile pinball), not just the mean; rewards forecasts both
  accurate \emph{and} honestly calibrated about uncertainty
  (probabilistic models only).
\end{itemize}

\textbf{Cov flag.} Each table carries a \textbf{Cov} column: \(\checkmark\) = the
model ingests the weather/solar covariate track

\subsubsection{2.1 Reference / statistical}\label{reference-statistical}

\begin{longtable}[]{@{}
  >{\raggedright\arraybackslash}p{(\linewidth - 12\tabcolsep) * \real{0.0500}}
  >{\raggedright\arraybackslash}p{(\linewidth - 12\tabcolsep) * \real{0.1500}}
  >{\raggedright\arraybackslash}p{(\linewidth - 12\tabcolsep) * \real{0.0700}}
  >{\raggedright\arraybackslash}p{(\linewidth - 12\tabcolsep) * \real{0.0700}}
  >{\raggedright\arraybackslash}p{(\linewidth - 12\tabcolsep) * \real{0.0800}}
  >{\raggedright\arraybackslash}p{(\linewidth - 12\tabcolsep) * \real{0.0500}}
  >{\raggedright\arraybackslash}p{(\linewidth - 12\tabcolsep) * \real{0.5300}}@{}}
\toprule\noalign{}
\begin{minipage}[b]{\linewidth}\raggedright
Rank
\end{minipage} & \begin{minipage}[b]{\linewidth}\raggedright
Model
\end{minipage} & \begin{minipage}[b]{\linewidth}\raggedright
\textbf{SS \(\uparrow\)}
\end{minipage} & \begin{minipage}[b]{\linewidth}\raggedright
R\textsuperscript{2} \(\uparrow\)
\end{minipage} & \begin{minipage}[b]{\linewidth}\raggedright
CRPS \(\downarrow\)
\end{minipage} & \begin{minipage}[b]{\linewidth}\raggedright
Cov
\end{minipage} & \begin{minipage}[b]{\linewidth}\raggedright
Description
\end{minipage} \\
\midrule\noalign{}
\endhead
\bottomrule\noalign{}
\endlastfoot
23 & climatology\_\allowbreak{}hourly & 0.234 & 0.314 & --- & \(\times\) & Per-hour-of-day
mean power from the training set; a fixed seasonal climatology. \\
\midrule\noalign{}
28 & seasonal\_\allowbreak{}naive & 0.107 & 0.276 & --- & \(\times\) & Copies power from the
same clock time on the previous day. \\
\midrule\noalign{}
30 & persistence & 0.014 & 0.171 & --- & \(\times\) & Carries the last observed
power flat across the horizon. \\
\midrule\noalign{}
--- & smart\_\allowbreak{}persistence & 0.000 & 0.210 & --- & \(\times\) & Persistence
rescaled by the clear-sky ratio --- the sun-adjusted reference floor \\
\end{longtable}

\newpage

\subsubsection{2.2 Classical ML}\label{classical-ml}

\begin{longtable}[]{@{}
  >{\raggedright\arraybackslash}p{(\linewidth - 12\tabcolsep) * \real{0.0500}}
  >{\raggedright\arraybackslash}p{(\linewidth - 12\tabcolsep) * \real{0.1500}}
  >{\raggedright\arraybackslash}p{(\linewidth - 12\tabcolsep) * \real{0.0700}}
  >{\raggedright\arraybackslash}p{(\linewidth - 12\tabcolsep) * \real{0.0700}}
  >{\raggedright\arraybackslash}p{(\linewidth - 12\tabcolsep) * \real{0.0800}}
  >{\raggedright\arraybackslash}p{(\linewidth - 12\tabcolsep) * \real{0.0500}}
  >{\raggedright\arraybackslash}p{(\linewidth - 12\tabcolsep) * \real{0.5300}}@{}}
\toprule\noalign{}
\begin{minipage}[b]{\linewidth}\raggedright
Rank
\end{minipage} & \begin{minipage}[b]{\linewidth}\raggedright
Model
\end{minipage} & \begin{minipage}[b]{\linewidth}\raggedright
\textbf{SS \(\uparrow\)}
\end{minipage} & \begin{minipage}[b]{\linewidth}\raggedright
R\textsuperscript{2} \(\uparrow\)
\end{minipage} & \begin{minipage}[b]{\linewidth}\raggedright
CRPS \(\downarrow\)
\end{minipage} & \begin{minipage}[b]{\linewidth}\raggedright
Cov
\end{minipage} & \begin{minipage}[b]{\linewidth}\raggedright
Description
\end{minipage} \\
\midrule\noalign{}
\endhead
\bottomrule\noalign{}
\endlastfoot
11 & LightGBM & 0.384 & 0.555 & 0.0768 & \(\checkmark\) & Gradient-boosted decision
trees over lagged power + covariate features. \\
\end{longtable}

\subsubsection{2.3 Tabular foundation
models}\label{tabular-foundation-models}

\begin{longtable}[]{@{}
  >{\raggedright\arraybackslash}p{(\linewidth - 12\tabcolsep) * \real{0.0500}}
  >{\raggedright\arraybackslash}p{(\linewidth - 12\tabcolsep) * \real{0.1500}}
  >{\raggedright\arraybackslash}p{(\linewidth - 12\tabcolsep) * \real{0.0700}}
  >{\raggedright\arraybackslash}p{(\linewidth - 12\tabcolsep) * \real{0.0700}}
  >{\raggedright\arraybackslash}p{(\linewidth - 12\tabcolsep) * \real{0.0800}}
  >{\raggedright\arraybackslash}p{(\linewidth - 12\tabcolsep) * \real{0.0500}}
  >{\raggedright\arraybackslash}p{(\linewidth - 12\tabcolsep) * \real{0.5300}}@{}}
\toprule\noalign{}
\begin{minipage}[b]{\linewidth}\raggedright
Rank
\end{minipage} & \begin{minipage}[b]{\linewidth}\raggedright
Model
\end{minipage} & \begin{minipage}[b]{\linewidth}\raggedright
\textbf{SS \(\uparrow\)}
\end{minipage} & \begin{minipage}[b]{\linewidth}\raggedright
R\textsuperscript{2} \(\uparrow\)
\end{minipage} & \begin{minipage}[b]{\linewidth}\raggedright
CRPS \(\downarrow\)
\end{minipage} & \begin{minipage}[b]{\linewidth}\raggedright
Cov
\end{minipage} & \begin{minipage}[b]{\linewidth}\raggedright
Paper
\end{minipage} \\
\midrule\noalign{}
\endhead
\bottomrule\noalign{}
\endlastfoot
17 & TabPFN 3 & 0.339 & 0.499 & 0.0815 & \(\checkmark\) &
\href{https://arxiv.org/abs/2605.13986}{arXiv:2605.13986} \\
\midrule\noalign{}
--- & TabFM 1.0 & TBD & --- & TBD & \(\checkmark\) &
\href{https://research.google/blog/introducing-tabfm-a-zero-shot-foundation-model-for-tabular-data/}{Google
Research} \\
\end{longtable}

\subsubsection{2.4 Supervised DL}\label{supervised-dl}

\begin{longtable}[]{@{}
  >{\raggedright\arraybackslash}p{(\linewidth - 12\tabcolsep) * \real{0.0500}}
  >{\raggedright\arraybackslash}p{(\linewidth - 12\tabcolsep) * \real{0.1500}}
  >{\raggedright\arraybackslash}p{(\linewidth - 12\tabcolsep) * \real{0.0700}}
  >{\raggedright\arraybackslash}p{(\linewidth - 12\tabcolsep) * \real{0.0700}}
  >{\raggedright\arraybackslash}p{(\linewidth - 12\tabcolsep) * \real{0.0800}}
  >{\raggedright\arraybackslash}p{(\linewidth - 12\tabcolsep) * \real{0.0500}}
  >{\raggedright\arraybackslash}p{(\linewidth - 12\tabcolsep) * \real{0.5300}}@{}}
\toprule\noalign{}
\begin{minipage}[b]{\linewidth}\raggedright
Rank
\end{minipage} & \begin{minipage}[b]{\linewidth}\raggedright
Model
\end{minipage} & \begin{minipage}[b]{\linewidth}\raggedright
\textbf{SS \(\uparrow\)}
\end{minipage} & \begin{minipage}[b]{\linewidth}\raggedright
R\textsuperscript{2} \(\uparrow\)
\end{minipage} & \begin{minipage}[b]{\linewidth}\raggedright
CRPS \(\downarrow\)
\end{minipage} & \begin{minipage}[b]{\linewidth}\raggedright
Cov
\end{minipage} & \begin{minipage}[b]{\linewidth}\raggedright
Description
\end{minipage} \\
\midrule\noalign{}
\endhead
\bottomrule\noalign{}
\endlastfoot
1 & \textbf{iTransformer} & \textbf{0.552} & 0.765 & --- & \(\checkmark\) &
\href{https://arxiv.org/pdf/2310.06625}{arXiv:2310.06625} \\
\midrule\noalign{}
7 & PatchTST & 0.458 & 0.651 & --- & \(\checkmark\) & Splits each series into patches
for a channel-independent transformer. \\
\midrule\noalign{}
9 & TFT & 0.423 & 0.626 & 0.0689 & \(\checkmark\) & Temporal Fusion Transformer;
gated attention with quantile outputs. \\
\midrule\noalign{}
10 & MLP & 0.413 & 0.600 & --- & \(\checkmark\) & Plain multilayer perceptron over
the flattened history window. \\
\midrule\noalign{}
20 & DLinear & 0.325 & 0.462 & --- & \(\times\) & Single linear layer on
trend/seasonal decomposition --- target-only simplicity check. \\
\end{longtable}

\subsubsection{2.5 Zero-shot TSFM}\label{zero-shot-tsfm}

\begin{longtable}[]{@{}
  >{\raggedright\arraybackslash}p{(\linewidth - 12\tabcolsep) * \real{0.0500}}
  >{\raggedright\arraybackslash}p{(\linewidth - 12\tabcolsep) * \real{0.1500}}
  >{\raggedright\arraybackslash}p{(\linewidth - 12\tabcolsep) * \real{0.0700}}
  >{\raggedright\arraybackslash}p{(\linewidth - 12\tabcolsep) * \real{0.0700}}
  >{\raggedright\arraybackslash}p{(\linewidth - 12\tabcolsep) * \real{0.0800}}
  >{\raggedright\arraybackslash}p{(\linewidth - 12\tabcolsep) * \real{0.0500}}
  >{\raggedright\arraybackslash}p{(\linewidth - 12\tabcolsep) * \real{0.5300}}@{}}
\toprule\noalign{}
\begin{minipage}[b]{\linewidth}\raggedright
Rank
\end{minipage} & \begin{minipage}[b]{\linewidth}\raggedright
Model
\end{minipage} & \begin{minipage}[b]{\linewidth}\raggedright
\textbf{SS \(\uparrow\)}
\end{minipage} & \begin{minipage}[b]{\linewidth}\raggedright
R\textsuperscript{2} \(\uparrow\)
\end{minipage} & \begin{minipage}[b]{\linewidth}\raggedright
CRPS \(\downarrow\)
\end{minipage} & \begin{minipage}[b]{\linewidth}\raggedright
Cov
\end{minipage} & \begin{minipage}[b]{\linewidth}\raggedright
Paper
\end{minipage} \\
\midrule\noalign{}
\endhead
\bottomrule\noalign{}
\endlastfoot
21 & tirex\_\allowbreak{}zs & 0.287 & 0.440 & 0.0892 & \(\times\) &
\href{https://arxiv.org/abs/2505.23719}{arXiv:2505.23719} \\
\midrule\noalign{}
22 & timesfm\_\allowbreak{}zs & 0.271 & 0.429 & 0.0923 & \(\times\) &
\href{https://huggingface.co/google/timesfm-2.5-200m-pytorch}{TimesFM
2.5} \\
\midrule\noalign{}
26 & chronos2\_\allowbreak{}zs & 0.187 & 0.277 & 0.1072 & \(\checkmark\) &
\href{https://arxiv.org/abs/2510.15821}{arXiv:2510.15821} \\
\midrule\noalign{}
\(\times\) & ttm\_\allowbreak{}zs & \textbf{\(-\)0.081} & 0.155 & --- & \(\checkmark\) &
\href{https://huggingface.co/ibm-granite/granite-timeseries-ttm-r2}{TTM-R2} \\
\end{longtable}

\subsubsection{2.6 Fine-tuned TSFM}\label{fine-tuned-tsfm}

\begin{longtable}[]{@{}
  >{\raggedright\arraybackslash}p{(\linewidth - 12\tabcolsep) * \real{0.0500}}
  >{\raggedright\arraybackslash}p{(\linewidth - 12\tabcolsep) * \real{0.1500}}
  >{\raggedright\arraybackslash}p{(\linewidth - 12\tabcolsep) * \real{0.0700}}
  >{\raggedright\arraybackslash}p{(\linewidth - 12\tabcolsep) * \real{0.0700}}
  >{\raggedright\arraybackslash}p{(\linewidth - 12\tabcolsep) * \real{0.0800}}
  >{\raggedright\arraybackslash}p{(\linewidth - 12\tabcolsep) * \real{0.0500}}
  >{\raggedright\arraybackslash}p{(\linewidth - 12\tabcolsep) * \real{0.5300}}@{}}
\toprule\noalign{}
\begin{minipage}[b]{\linewidth}\raggedright
Rank
\end{minipage} & \begin{minipage}[b]{\linewidth}\raggedright
Model
\end{minipage} & \begin{minipage}[b]{\linewidth}\raggedright
\textbf{SS \(\uparrow\)}
\end{minipage} & \begin{minipage}[b]{\linewidth}\raggedright
R\textsuperscript{2} \(\uparrow\)
\end{minipage} & \begin{minipage}[b]{\linewidth}\raggedright
CRPS \(\downarrow\)
\end{minipage} & \begin{minipage}[b]{\linewidth}\raggedright
Cov
\end{minipage} & \begin{minipage}[b]{\linewidth}\raggedright
Paper
\end{minipage} \\
\midrule\noalign{}
\endhead
\bottomrule\noalign{}
\endlastfoot
3 & chronos2\_\allowbreak{}oracle\_\allowbreak{}ft & 0.504 & 0.706 & 0.0630 & \(\checkmark\) &
\href{https://arxiv.org/abs/2510.15821}{arXiv:2510.15821} \\
\midrule\noalign{}
6 & chronos2\_\allowbreak{}oracle & 0.474 & 0.679 & 0.0635 & \(\checkmark\) &
\href{https://arxiv.org/abs/2510.15821}{arXiv:2510.15821} \\
\midrule\noalign{}
13 & ttm\_\allowbreak{}ft & 0.364 & 0.520 & --- & \(\checkmark\) &
\href{https://huggingface.co/ibm-granite/granite-timeseries-ttm-r2}{TTM-R2} \\
\midrule\noalign{}
19 & chronos2\_\allowbreak{}ft & 0.331 & 0.475 & 0.0855 & \(\checkmark\) &
\href{https://arxiv.org/abs/2510.15821}{arXiv:2510.15821} \\
\end{longtable}

\begin{quote}
\textbf{\texttt{chronos2\_\allowbreak{}*} vs \texttt{chronos2\_\allowbreak{}oracle\_\allowbreak{}*}} --- same
Chronos-2 model; the only difference is \emph{which covariates are
exposed over the forecast horizon}. The plain variants
(\texttt{chronos2\_\allowbreak{}zs}, \texttt{chronos2\_\allowbreak{}ft}) uses only deterministic
future covariates (solar geometry + calendar) that are known ahead. The
\texttt{oracle} variants feed the model the \textbf{future observed
weather} (temperature, cloud cover, irradiance, \ldots) over the
horizon.
\end{quote}

\subsubsection{2.7 Retrieval / frozen-FM
adaptation}\label{retrieval-frozen-fm-adaptation}

\begin{longtable}[]{@{}
  >{\raggedright\arraybackslash}p{(\linewidth - 14\tabcolsep) * \real{0.0500}}
  >{\raggedright\arraybackslash}p{(\linewidth - 14\tabcolsep) * \real{0.1200}}
  >{\raggedright\arraybackslash}p{(\linewidth - 14\tabcolsep) * \real{0.0700}}
  >{\raggedright\arraybackslash}p{(\linewidth - 14\tabcolsep) * \real{0.0700}}
  >{\raggedright\arraybackslash}p{(\linewidth - 14\tabcolsep) * \real{0.0800}}
  >{\raggedright\arraybackslash}p{(\linewidth - 14\tabcolsep) * \real{0.0500}}
  >{\raggedright\arraybackslash}p{(\linewidth - 14\tabcolsep) * \real{0.2800}}
  >{\raggedright\arraybackslash}p{(\linewidth - 14\tabcolsep) * \real{0.2800}}@{}}
\toprule\noalign{}
\begin{minipage}[b]{\linewidth}\raggedright
Rank
\end{minipage} & \begin{minipage}[b]{\linewidth}\raggedright
Model
\end{minipage} & \begin{minipage}[b]{\linewidth}\raggedright
\textbf{SS \(\uparrow\)}
\end{minipage} & \begin{minipage}[b]{\linewidth}\raggedright
R\textsuperscript{2} \(\uparrow\)
\end{minipage} & \begin{minipage}[b]{\linewidth}\raggedright
CRPS \(\downarrow\)
\end{minipage} & \begin{minipage}[b]{\linewidth}\raggedright
Cov
\end{minipage} & \begin{minipage}[b]{\linewidth}\raggedright
Regime
\end{minipage} & \begin{minipage}[b]{\linewidth}\raggedright
Paper
\end{minipage} \\
\midrule\noalign{}
\endhead
\bottomrule\noalign{}
\endlastfoot
4 & ts\_\allowbreak{}rag & 0.478 & 0.331 & --- & \(\times\) & Zero-shot (frozen FM +
retrieval) &
\href{https://arxiv.org/abs/2503.07649}{arXiv:2503.07649} \\
\midrule\noalign{}
5 & cross\_\allowbreak{}rag & 0.477 & 0.375 & --- & \(\times\) & Zero-shot (frozen FM +
retrieval) &
\href{https://arxiv.org/abs/2603.14709}{arXiv:2603.14709} \\
\midrule\noalign{}
12 & cora & 0.374 & 0.550 & 0.0816 & \(\checkmark\) & Fine-tuned (adapter, frozen FM)
& \href{https://arxiv.org/abs/2510.12681}{arXiv:2510.12681} \\
\end{longtable}

\subsubsection{2.8 Endogenous multimodal (second modality synthesized
from the numeric track --- no external
sensor)}\label{endogenous-multimodal-second-modality-synthesized-from-the-numeric-track-no-external-sensor}

Pseudo-images rendered from the series (time\_\allowbreak{}vlm, visionts\_\allowbreak{}pp) or
weather-text templated from the covariates (aurora). No satellite/sky
frames consumed.

\begin{longtable}[]{@{}
  >{\raggedright\arraybackslash}p{(\linewidth - 14\tabcolsep) * \real{0.0500}}
  >{\raggedright\arraybackslash}p{(\linewidth - 14\tabcolsep) * \real{0.1200}}
  >{\raggedright\arraybackslash}p{(\linewidth - 14\tabcolsep) * \real{0.0700}}
  >{\raggedright\arraybackslash}p{(\linewidth - 14\tabcolsep) * \real{0.0700}}
  >{\raggedright\arraybackslash}p{(\linewidth - 14\tabcolsep) * \real{0.0800}}
  >{\raggedright\arraybackslash}p{(\linewidth - 14\tabcolsep) * \real{0.0500}}
  >{\raggedright\arraybackslash}p{(\linewidth - 14\tabcolsep) * \real{0.2800}}
  >{\raggedright\arraybackslash}p{(\linewidth - 14\tabcolsep) * \real{0.2800}}@{}}
\toprule\noalign{}
\begin{minipage}[b]{\linewidth}\raggedright
Rank
\end{minipage} & \begin{minipage}[b]{\linewidth}\raggedright
Model
\end{minipage} & \begin{minipage}[b]{\linewidth}\raggedright
\textbf{SS \(\uparrow\)}
\end{minipage} & \begin{minipage}[b]{\linewidth}\raggedright
R\textsuperscript{2} \(\uparrow\)
\end{minipage} & \begin{minipage}[b]{\linewidth}\raggedright
CRPS \(\downarrow\)
\end{minipage} & \begin{minipage}[b]{\linewidth}\raggedright
Cov
\end{minipage} & \begin{minipage}[b]{\linewidth}\raggedright
Regime
\end{minipage} & \begin{minipage}[b]{\linewidth}\raggedright
Paper
\end{minipage} \\
\midrule\noalign{}
\endhead
\bottomrule\noalign{}
\endlastfoot
2 & \textbf{time\_\allowbreak{}vlm} & \textbf{0.540} & 0.487 & --- & \(\times\) & Fine-tuned
(frozen VLM backbones) &
\href{https://arxiv.org/abs/2502.04395}{arXiv:2502.04395} \\
\midrule\noalign{}
24 & aurora & 0.232 & 0.440 & --- & \(\checkmark\) & Zero-shot &
\href{https://arxiv.org/abs/2509.22295}{arXiv:2509.22295} \\
\midrule\noalign{}
29 & visionts\_\allowbreak{}pp & 0.017 & 0.397 & --- & \(\times\) & Zero-shot (frozen image
MAE) & \href{https://arxiv.org/abs/2508.04379}{arXiv:2508.04379} \\
\end{longtable}

\subsubsection{2.9 Exogenous multimodal (external satellite
imagery)}\label{exogenous-multimodal-external-satellite-imagery}

\begin{longtable}[]{@{}
  >{\raggedright\arraybackslash}p{(\linewidth - 14\tabcolsep) * \real{0.0500}}
  >{\raggedright\arraybackslash}p{(\linewidth - 14\tabcolsep) * \real{0.1200}}
  >{\raggedright\arraybackslash}p{(\linewidth - 14\tabcolsep) * \real{0.0700}}
  >{\raggedright\arraybackslash}p{(\linewidth - 14\tabcolsep) * \real{0.0700}}
  >{\raggedright\arraybackslash}p{(\linewidth - 14\tabcolsep) * \real{0.0800}}
  >{\raggedright\arraybackslash}p{(\linewidth - 14\tabcolsep) * \real{0.0500}}
  >{\raggedright\arraybackslash}p{(\linewidth - 14\tabcolsep) * \real{0.2800}}
  >{\raggedright\arraybackslash}p{(\linewidth - 14\tabcolsep) * \real{0.2800}}@{}}
\toprule\noalign{}
\begin{minipage}[b]{\linewidth}\raggedright
Rank
\end{minipage} & \begin{minipage}[b]{\linewidth}\raggedright
Model
\end{minipage} & \begin{minipage}[b]{\linewidth}\raggedright
\textbf{SS \(\uparrow\)}
\end{minipage} & \begin{minipage}[b]{\linewidth}\raggedright
R\textsuperscript{2} \(\uparrow\)
\end{minipage} & \begin{minipage}[b]{\linewidth}\raggedright
CRPS \(\downarrow\)
\end{minipage} & \begin{minipage}[b]{\linewidth}\raggedright
Cov
\end{minipage} & \begin{minipage}[b]{\linewidth}\raggedright
Regime
\end{minipage} & \begin{minipage}[b]{\linewidth}\raggedright
Paper
\end{minipage} \\
\midrule\noalign{}
\endhead
\bottomrule\noalign{}
\endlastfoot
8 & solar\_\allowbreak{}vlm & 0.443 & 0.660 & --- & \(\checkmark\) & Fine-tuned (pretrained
vision) & \href{https://arxiv.org/abs/2604.04145}{arXiv:2604.04145} \\
\midrule\noalign{}
14 & crossvivit & 0.349 & 0.565 & --- & \(\checkmark\) & From scratch &
\href{https://arxiv.org/abs/2306.01112}{arXiv:2306.01112} \\
\midrule\noalign{}
16 & \textbf{mmtsfm\_\allowbreak{}grassmann\_\allowbreak{}interleaved} & \textbf{0.3429} & --- &
0.0805 & \(\checkmark\) & Fine-tuned (frozen FM + vision) & ours (this work) \\
\midrule\noalign{}
25 & sunset & 0.216 & 0.359 & --- & \(\times\) & From scratch &
\href{https://github.com/YuchiSun/SUNSET}{github: YuchiSun/SUNSET} \\
\midrule\noalign{}
27 & unicast & 0.121 & 0.371 & --- & \(\times\) & Fine-tuned (soft-prompt, frozen
FM) & \href{https://arxiv.org/abs/2508.11954}{arXiv:2508.11954} \\
\midrule\noalign{}
--- & FusionSF & TBD & --- & TBD & \(\checkmark\) & TBD &
\href{https://arxiv.org/abs/2402.05823}{arXiv:2402.05823} \\
\midrule\noalign{}
--- & PIPE & TBD & --- & TBD & \(\checkmark\) & TBD &
\href{https://arxiv.org/abs/2506.14786}{arXiv:2506.14786} \\
\midrule\noalign{}
--- & MATE & TBD & --- & TBD & \(\checkmark\) & TBD &
\href{https://openreview.net/forum?id=jn7GJdyuVn}{OpenReview} \\
\end{longtable}

\newpage

\subsection{4. Conclusions}\label{conclusions}

Drawn from the cross-plant leaderboard. All numbers are zero-shot on
\textbf{disjoint test plants} --- no test-plant data in training.

\textbf{1. A well-tuned supervised transformer is the model to beat.}
iTransformer tops every metric (SS 0.552, R\textsuperscript{2} 0.765), ahead of every
foundation model, retrieval scheme, and multimodal system. On this
cross-plant PV task, architecture + in-domain supervised training still
outperforms large pretrained backbones.

\textbf{2. Zero-shot time-series foundation models underdeliver on PV.}
The zero-shot TSFMs (chronos2\_\allowbreak{}zs SS 0.187, timesfm\_\allowbreak{}zs 0.271, tirex\_\allowbreak{}zs
0.287, ttm\_\allowbreak{}zs \textbf{\(-\)0.081} --- worse than the naive floor) trail
even classical baselines. PV's sharp, weather-driven dynamics are
out-of-distribution for generic TS pretraining; the models need
adaptation to be useful.

\textbf{3. Adaptation closes the gap, and retrieval is the best
frozen-backbone lever.} Every adaptation step helps: Chronos-2 jumps
zero-shot \(\rightarrow\) fine-tuned, and TTM flips from below-baseline to SS 0.364
once fine-tuned. Notably, \textbf{retrieval over a frozen backbone}
(ts\_\allowbreak{}rag SS 0.478, cross\_\allowbreak{}rag 0.477) rivals fine-tuning \textbf{without
updating any weights}, making it the most effective --- and cheapest ---
frozen-FM strategy tested.

\textbf{4. Multimodality has not yet paid off} The one multimodal model
at the top (time\_\allowbreak{}vlm, SS 0.540, rank 2) uses \emph{endogenous
pseudo-images synthesized from the numeric series} and no covariates.
Every model that ingests genuine \textbf{exogenous satellite/sky
imagery} ranks mid-pack or worse.

\textbf{7. R\textsuperscript{2} and SS disagree for biased forecasters --- read them
together.} Several models track the day's \emph{shape} well yet sit off
the diagonal: e.g.~visionts\_\allowbreak{}pp has R\textsuperscript{2} 0.397 but SS 0.017, and
solar\_\allowbreak{}vlm / sunset / crossvivit show systematic under-prediction. High
variance-tracking (R\textsuperscript{2}) with low skill (SS) flags a \textbf{bias/scale}
problem, not a timing problem --- fixable by calibration rather than a
better temporal model.

\end{document}
